% ATTEMPT_STATUS: unresolved
% RESEARCH_GENERATED: 2026-10-01; GPT-6 Astra Ultra; individual mathematical attempt
% TOP_PROBLEM: {"id": 102, "title": "Tuples in Dense Sets", "category_id": 2, "category": {"id": 2, "name": "combinatorics", "display_name": "Combinatorics", "description": "Counting problems, graph theory, discrete structures.", "slug": "combinatorics", "order_index": 2, "created_at": "2026-07-31T15:26:25.671Z"}, "status": "open"}
% FIRST_POSTED: 2026-10-01
% Imported from: unsolvedmath_additional_500.tex; UnsolvedMath ID 102
% !TeX program = lualatex
% Compile twice: lualatex 102.tex
% Self-contained: no external bibliography, images, or \input files.
% Numeric section labels and visible numbers are original UnsolvedMath IDs.
\documentclass[11pt,a4paper]{article}
\usepackage[margin=24mm,headheight=14pt]{geometry}
\usepackage{fontspec}
\setmainfont{Latin Modern Roman}
\setsansfont{Latin Modern Sans}
\setmonofont{Latin Modern Mono}
\usepackage{amsmath,amssymb,mathtools}
\usepackage{unicode-math}
\setmathfont{Latin Modern Math}
\usepackage{microtype}
\usepackage{enumitem,xurl,needspace}
\usepackage[dvipsnames]{xcolor}
\usepackage{hyperref}
\hypersetup{unicode=true,colorlinks=true,linkcolor=MidnightBlue,urlcolor=MidnightBlue,
 pdftitle={UnsolvedMath Problem 102},pdfauthor={Source-annotated compilation},bookmarksnumbered=true}
\usepackage{bookmark,tocloft,titlesec,fancyhdr}
\setlength{\cftsecnumwidth}{4.2em}
\cftsetpnumwidth{2.5em}
\cftsetrmarg{4em}
\titleformat{\section}{\Large\bfseries\raggedright}{\thesection}{0.7em}{}
\pagestyle{fancy}\fancyhf{}
\fancyhead[L]{\small\sffamily UnsolvedMath research attempt}
\fancyhead[R]{\small Original catalogue IDs}
\fancyfoot[C]{\thepage}
\setlength{\parindent}{0pt}
\setlength{\parskip}{5pt plus 1pt}
\setlength{\emergencystretch}{3em}
\setcounter{tocdepth}{1}\setcounter{secnumdepth}{1}
\setlist[itemize]{leftmargin=3em,itemsep=4pt,topsep=4pt,parsep=0pt}
\allowdisplaybreaks
\newcommand{\N}{\mathbb{N}}
\newcommand{\Z}{\mathbb{Z}}
\newcommand{\Q}{\mathbb{Q}}
\newcommand{\R}{\mathbb{R}}
\newcommand{\C}{\mathbb{C}}
\newcommand{\F}{\mathbb{F}}
\newcommand{\A}{\mathbb{A}}
\newcommand{\PP}{\mathbb{P}}
\newcommand{\E}{\mathbb{E}}
\newcommand{\eps}{\varepsilon}
\newcommand{\dd}{\,\mathrm d}
\newcommand{\sref}[2]{\hyperlink{source:#1:#2}{[S#2]}}
\newif\ifshowcatalogue
\showcataloguetrue
\newcommand{\cataloguescope}[1]{\ifshowcatalogue
 \paragraph{Statement recorded in the supplied source.}
 #1\fi}
\begin{document}
% UnsolvedMath ID 102; source catalogue code GREEN-012

\setcounter{section}{101}

\section{A Fifteen-Constraint Additive Pattern}\label{102}

{\small\sffamily UnsolvedMath ID 102 · Combinatorics}

\subsection{Definitions and mathematical statement}

\paragraph{Shared notation.}Unless a section says otherwise, $\N=\{1,2,\ldots\}$,
$\Z$ is the ring of integers, $\Q,\R,\C$ are the rational, real, and complex
fields, $[N]=\{1,\ldots,\lfloor N\rfloor\}$, and $\log$ is the natural
logarithm. For real functions of a parameter tending to infinity,
$f\ll g$ and $f=O(g)$ mean $|f|\leq Cg$ eventually for some fixed $C$;
$f\gg g$ reverses this comparison; $f=o(g)$ means $f/g\to0$;
$f\sim g$ means $f/g\to1$; and $f\asymp g$ means both order bounds.
A subscript on $O$ or $\ll$ specifies allowed constant dependence.
The expression $n^{o(1)}$ in an upper bound means growth smaller than
$n^\epsilon$ for each fixed $\epsilon>0$. Local definitions govern any
exception, finite-field convention, or use of the same symbol for another
object. An asymptotic claim is not automatically an effective algorithm.



All indices $i,j$ below are read modulo 5, in $\{1,\ldots,5\}$. The density is $\alpha=|A|/N$. Count ordered ten-tuples, with coincidences between entries permitted, through
\[T(A)=\sum_{x_1,\ldots,x_5,y_1,\ldots,y_5\in G}\prod_{i=1}^5\prod_{s=0}^2 1_A(x_i+y_{i+s}).\]
The precise proposed inequality is $T(A)\geq\alpha^{15}N^{10}$ for every finite abelian group and every subset. \sref{102}{1} \sref{102}{2}

\paragraph{Statement recorded in the supplied source.}
Let $G$ be an abelian group of size $N$, and suppose that $A \subset G$ has density $\alpha$. Are there at least $\alpha^{15}N^{10}$ tuples $(x_1, \ldots, x_5, y_1, \ldots, y_5) \in G^{10}$ such that $x_i + y_j \in A$ whenever $j \in \{i, i+1, i+2\}$? \sref{102}{1}

\paragraph{Residual task reported by the archive.}

The embedded review (2026-08-17) records: Produce a counterexample or prove the stated lower bound. \sref{102}{1}

\subsection{Short English statement}

Does every dense subset of a finite abelian group contain at least the random-model number of configurations satisfying these fifteen linked sum conditions?

\subsection{Research attempt}

\paragraph{Provenance and scope.}
This individual attempt was generated by GPT-6 Astra Ultra on 1 October 2026. The method was to derive explicit partial results and test the first proposed extension adversarially. The original statement and sources below are retained. No claim of mathematical novelty or complete solution is made.

\paragraph{Formal target and source check.}
Let $f=1_A$ on a finite abelian group $G$ of order $N$, and let $\alpha=|A|/N$. The desired normalised count is $t(A)=N^{-10}T(A)\geq\alpha^{15}$. All ten variables are ordered and may coincide. The author-hosted list identifies the relevant graph as the bipartite graph $K_{5,5}$ with a ten-cycle removed. No general Sidorenko theorem for that graph is assumed in the following work. We attack the special translation-invariant kernel using exact elimination of five variables, then test structured sets and perturbations.

\paragraph{An exact five-variable reduction.}
For $u,v,w\in G$, define
\[
 C(u,v,w)=\E_x f(x+u)f(x+v)f(x+w).
\]
Conditioning on $y_1,\ldots,y_5$, the five $x_i$ are independent, so
\[
 t(A)=\E_{y_1,\ldots,y_5}\prod_{i=1}^5
 C(y_i,y_{i+1},y_{i+2}),
\]
where indices are cyclic. Also $\E_{u,v,w}C(u,v,w)=\alpha^3$, because after fixing $x$, the three translates are independent uniform points. Thus the target would follow from a positive-correlation inequality for these five overlapping triple correlations. Jensen cannot be applied directly to replace each factor by its mean: the factors are different functions of overlapping variables, and a product of nonnegative functions can have expectation below the product of expectations. This is the exact missing inequality for this route, not a proved intermediate step.

\paragraph{Cosets satisfy a stronger bound exactly.}
Suppose $A=a+H$ for a subgroup $H$ of index $q$, so $\alpha=1/q$. In the quotient $G/H$, every required edge has the equation $\bar x_i+\bar y_j=\bar a$. The bipartite graph is connected: neighbouring $x$-vertices share a $y$-neighbour, hence all $\bar x_i$ are equal and all $\bar y_j$ are the complementary value. There are exactly $q$ assignments in the quotient out of $q^{10}$ possibilities. Each lifts to $|H|^{10}$ assignments. Therefore
\[
 T(A)=N^{10}q^{-9}=\alpha^9N^{10}\geq\alpha^{15}N^{10}.
\]
This calculation handles the full group and a singleton as endpoint cases. It also shows that subgroup examples have an excess by a factor $\alpha^{-6}$, so they are poor candidates for a counterexample when the density is small.

\paragraph{Product sets amplify any real counterexample.}
For $A\subset G$ and $B\subset H$, each of the fifteen conditions for $A\times B$ splits into an independent condition in the two coordinates. Consequently
\[
 t(A\times B)=t(A)t(B),\quad
 \frac{t(A\times B)}{(\alpha\beta)^{15}}
 =\frac{t(A)}{\alpha^{15}}\frac{t(B)}{\beta^{15}}.
\]
Thus even a small strict counterexample in one finite abelian group would yield arbitrarily strong multiplicative deficits by taking powers. Conversely, successful finite checks do not prove the inequality in larger groups: they only eliminate those base examples.

\paragraph{Perturbing the constant kernel: the first surviving term.}
Consider the bounded real kernel $f=\alpha+\varepsilon g$, with $\E g=0$. Expand the fifteen-edge product. Any chosen collection of $g$-edges with a vertex of degree one contributes zero, since averaging its vertex variable averages a single translate of $g$. In a simple bipartite graph, a nonempty edge set of minimum degree at least two has at least four edges; at four edges it must be a four-cycle. Hence all coefficients of orders one, two and three vanish. Each order-four term is
\[
 \E_{x,x',y,y'}g(x+y)g(x+y')g(x'+y)g(x'+y')
 =\E_{y,y'}\bigl(\E_x g(x+y)g(x+y')\bigr)^2\geq0.
\]
This proves that the first possible nonzero perturbative coefficient is nonnegative. It is an exact local calculation, not a proof that all higher coefficients are nonnegative. In particular a zero fourth-order coefficient, or a finite perturbation, cannot be decided from this expansion alone. The averaging square also requires real $g$; no unqualified positivity claim for arbitrary complex kernels is being made.

\paragraph{Verification and final gap.}
The companion script computes the exact count by the five-variable formula for every nonempty subset of cyclic groups of orders two, three, four and five. It compares integers $T(A)N^5$ and $|A|^{15}$, avoiding rounding. The output records no violations in those finite cases. The work establishes exact coset counts, multiplicativity, and nonnegativity of the leading perturbative term. It leaves unresolved the five-overlapping-correlations inequality for general indicator functions; higher-order correlations can neither be discarded nor assigned a sign by the argument above. The original universal lower bound is therefore not claimed.

\subsection{Sources}

{\small

\begin{itemize}

\item[\hypertarget{source:102:1}{S1}] ULAM.AI, \emph{UnsolvedMath}, supplied \texttt{problems.json}, record 102 (GREEN-012), statement and background. The embedded literature-review date is 2026-08-17; that is the archive’s review date, not a fresh certification. Dataset landing page: \url{https://huggingface.co/datasets/ulamai/UnsolvedMath}.

\item[\hypertarget{source:102:2}{S2}] Ben Green, 100 Open Problems (author’s maintained problem list). \url{https://people.maths.ox.ac.uk/greenbj/papers/open-problems.pdf}. The author-hosted document was inspected on 27 September 2026; the archive’s local code is not a problem-number locator in this paper.

\item[\hypertarget{source:102:3}{S3}] M. Deng, J. Tidor, Y. Zhao, Uniform sets with few progressions via colourings, Math. Proc. Camb. Phil. Soc. 179 (2025), 79--103. \url{https://doi.org/10.1017/S0305004125000106}. Bibliographic information imported from the supplied record.

\end{itemize}

}

\label{end:102}
\end{document}
